\documentclass[10pt,a4paper]{amsart}
\usepackage[margin=19mm]{geometry}
\usepackage{amsmath,amssymb,mathtools}
\usepackage[scr=rsfs,cal=euler]{mathalfa}
\usepackage{xfrac}
\usepackage[unicode,hidelinks,bookmarks=false]{hyperref}
\newcommand{\ie}{i.e.\ }
\newcommand{\cf}{cf.\ }
\newcommand{\resp}{respectively\ }
\newcommand{\Subsection}{\subsection}
\providecommand{\nobreakdash}{\nobreak\hbox{-}}
\setlength{\emergencystretch}{2em}
\multlinegap0pt
\usepackage{amssymb}
\usepackage{tikz}
\usepackage{float}
\usepackage[shortlabels]{enumitem}
\usepackage{mathtools}
\usepackage{breqn}
\newtheorem{thm}{Theorem}
\newcommand{\f}{\Bbb F}
\title{Classification of Rational Functions of Degree Three over Finite Fields}
\author{Xiang-dong Hou, Siyu Peng, Yongyu Qiang, Shujun Zhao}
\date{Source: arXiv:2604.18954v2 (CC BY 4.0)}
\begin{document}
\maketitle

\noindent\emph{Source and licence.} Classification of Rational Functions of Degree Three over Finite Fields. Version arXiv:2604.18954v2 (CC BY 4.0), distributed by arXiv under the Creative Commons Attribution 4.0 International licence, http://creativecommons.org/licenses/by/4.0/, as recorded in the arXiv licence metadata for this exact version and shown on the abstract page https://arxiv.org/abs/2604.18954v2. Full licence notice: \texttt{licenses/arxiv-2604.18954-CC-BY-4.0.txt}.\par\smallskip

\noindent\textit{Editorial scope.} Counted unit: the article's thm thm1 (source0.tex lines 1467--1476). The article's complete original proof of it is delivered with it (source0.tex lines 1478--1682, one \texttt{proof} environment of 205 lines). No mathematics is written, repaired or re-derived: the statement and the proof are the article's own lines, in the article's own notation, and no retained line is substituted or re-flowed.  No further source-local context is needed: every cross-reference of every retained line resolves inside the two delivered spans.  Nothing was omitted from any retained span, so the package carries no omission record.  No retained line contains a citation, so the wrapper carries no bibliography and no bibliography entry.  No retained line contains a figure environment, an image-inclusion command or drawing code, so no image is delivered and the package declares no asset dependency.  Editorial formatting only, in addition: an AMS \texttt{amsart} preamble that restates the article's own declarations and macros used by the retained lines, namely the environments \texttt{thm} on separate counters, together with the standard packages those lines need.  The retained environments print in this document as Theorem 1 in order of appearance, whereas the article prints the same items with the numbering of its own full document.  The complete native source of the article as distributed in the arXiv e-print for this exact version is bundled unchanged as \texttt{sources/198832/source0.tex}.
\begin{thm}\label{thm1}
Assume $p\ge 5$. We have 
\[
\Theta_{\rm b}=27+\mathcal A,
\]
where
\[
\mathcal A=\bigl\{(\mu(\mu^2-9)/(\mu^2-1))^2: \mu\in\f_q\setminus\{0,\pm1,\pm3\}\bigr\}.
\]
\end{thm}

\begin{proof}
We assume $q\ge 71$. For $q\le 67$, the claim is verified by computer. We fix a nonsquare $d$ in $\f_q$.

\medskip
Let $\alpha=a^2\in\Theta_{\rm b}-27$, where $a\in\f_q^*$. Then there exists $(s,t)\in\f_q^2$ with $t(1+s+t)\ne 0$ such that
\begin{equation}\label{27+alpha}
\frac{s^3}{t(1+s+t)}=27+\alpha
\end{equation}
and
\begin{equation}\label{Gsta0a1}
G_{s,t}=X^4-2X^3-sX^2-2tX+t=(X^2+a_1X+a_0)(X^2+b_1X+b_0),
\end{equation}
where $a_0,a_1,b_0,b_1\in\f_q$, and
\begin{equation}\label{du2}
a_1^2-4a_0=du^2
\end{equation}
for some $u\in\f_q^*$. It follows from \eqref{Gsta0a1} that
\begin{equation}\label{s=t=}
\begin{cases}
s=-a_0-b_0-a_1b_1,\cr
t=a_0b_0,
\end{cases}
\end{equation}
and
\[
\begin{cases}
a_1+b_1=-2,\cr
a_1b_0+a_0b_1=-2a_0b_0.
\end{cases}
\]
Thus $b_0$ and $b_1$ can be expressed in terms of $a_0$ and $a_1$:
\begin{equation}\label{b0b1}
\begin{cases}
b_0=\displaystyle\frac{a_0(2+a_1)}{2a_0+a_1},\vspace{0.4em}\cr
b_1=-2-a_1.
\end{cases}
\end{equation}
Note that $2a_0+a_1\ne 0$. (Otherwise, we have $2a_0+a_1=0$ and $a_0(2+a_1)=0$. It follows that $(a_0,a_1)=(0,0)$ or $(1,-2)$. In either case, $X^2+a_1X+a_0$ is not irreducible over $\f_q$.)
By \eqref{27+alpha}, \eqref{s=t=} and \eqref{b0b1}, we have
\begin{align}\label{a^2}
a^2\,&=\alpha=\frac{s^3}{t(1+s+t)}-27\\
&=-\frac{(2+2 a_0+a_1) (4 a_0-a_1^2) (a_0+a_0^2+5 a_0 a_1+2 a_1^2+2 a_0 a_1^2+a_1^3)^2}{a_0^2 a_1 (2+a_1) (1+a_0+a_1)^2 (2 a_0+a_1)}. \nonumber
\end{align}
Since $4a_0-a_1^2=-du^2$, it follows from \eqref{a^2} that 
\begin{equation}\label{cod2}
\frac{d(2+2 a_0+a_1)}{a_1 (2+a_1)(2 a_0+a_1)}=e^2
\end{equation}
for some $e\in\f_q^*$. By choosing a suitable sign for $a$, we have
\begin{equation}\label{cod1}
a=\frac{eu(a_0+a_0^2+5 a_0 a_1+2 a_1^2+2 a_0 a_1^2+a_1^3)}{a_0(1+a_0+a_1)}.
\end{equation}
Making the substitution $a_0=(a_1^2-du^2)/4$ in \eqref{cod2} and \eqref{cod1}, we arrive at 
\[
C_1(a_1)=0,
\]
\[
C_2(a_1)=0,
\]
where
\begin{equation}\label{C1}
C_1(X)=\frac 12 (4 d+2 X d+X^2 d-4 X^2 e^2-4 X^3 e^2-X^4 e^2-d^2 u^2+2 X d e^2 u^2+X^2 d e^2 u^2),
\end{equation}
\begin{align}\label{C2}
C_2(X)=\,&\frac1{16} (-4 a X^2 - 4 a X^3 - a X^4 + 36 X^2 e u + 36 X^3 e u + 
   9 X^4 e u + 4 a d u^2 \\
   &+ 4 a X d u^2 + 2 a X^2 d u^2 - 
   4 d e u^3 - 20 X d e u^3 - 10 X^2 d e u^3 - a d^2 u^4 + 
   d^2 e u^5). \nonumber
\end{align}
We find that
\[
\text{Res}(C_1,C_2;X)=16 d^4 (-4+d u^2)^2 (-a+9 e u+a e^2 u^2-e^3 u^3)^4.
\]
Hence
\[
-a+9 e u+a e^2 u^2-e^3 u^3=0,
\]
i.e.,
\[
a=\frac{eu((eu)^2-9)}{(eu)^2-1}.
\]
Therefore, $\alpha=a^2\in\mathcal A$.

\medskip
On the other hand, assume $\alpha\in\mathcal A$, say $\alpha=a^2$, where
\[
a=\frac{\mu(\mu^2-9)}{\mu^2-1}
\]
for some $\mu\in\f_q\setminus\{0,\pm1,\pm3\}$. Write $\mu=eu$, where $e,u\in\f_q^*$. We find that $C_1(X)$ and $C_2(X)$ defined in \eqref{C1} and \eqref{C2} are related by 
\[
C_2(X)=\frac{-eu^3}{e^2u^2-1}C_1(X).
\]
Using the relation $\mu=eu$, we may express $C_1(X)$ in terms of $\mu$ and $u$:
\[
C_1(X)=\frac {-1}{2u^2}D(X,u),
\]
where 
\begin{align}\label{DXY}
D(X,Y)=\,&4 \mu^2 X^2 + 4 \mu^2 X^3 + \mu^2 X^4 - 4 d Y^2 - 2 d X Y^2\\
& -2 d \mu^2 X Y^2 - d X^2 Y^2 - d \mu^2 X^2 Y^2 + d^2 Y^4.\nonumber
\end{align}

We claim that $D(X,Y)$ is absolutely irreducible. Write $D(X,Y)$ in terms of homogeneous parts:
\[
D=D_2+D_3+D_4,
\]
where
\begin{align*}
&D_2=4 (\mu^2 X^2-d Y^2),\cr
&D_3=2 X (2 \mu^2 X^2-d Y^2-d \mu^2 Y^2),\cr
&D_4=(X^2-d Y^2) (\mu^2 X^2-d Y^2).
\end{align*}
Since $\mu^2\ne 0,1$, we have
\begin{equation}\label{gcd=1}
\gcd(D_2,D_3)=1.
\end{equation}
Assume to the contrary that $D$ is reducible over $\f_q$. Because of \eqref{gcd=1},
\[
D=(A_0+A_1)(B_2+B_3)\quad\text{or}\quad D=(A_1+A_2)(B_1+B_2),
\]
where $A_i,B_j\in\overline\f_q[X,Y]$ are homogeneous of degree $i$ and $j$, respectively. 

\medskip
{\bf Case 1.} Assume $D=(A_0+A_1)(B_2+B_3)$.

We may assume $A_0=1$ and $B_2=D_2$. Since $A_1B_3=D_4$, we have $B_3\mid D_4$. It follows that $\gcd(B_2,B_3)\ne 1$, which contradicts \eqref{gcd=1}.

\medskip
{\bf Case 2.} Assume $D=(A_1+A_2)(B_1+B_2)$.  

Since $A_1B_1=D_2=4 (\mu^2 X^2-d Y^2)$, we may assume that 
\begin{align*}
&A_1=\mu X-\delta Y,\cr
&B_1=4(\mu X+\delta Y),
\end{align*}
where $\delta\in\f_{q^2}\setminus\f_q$ is such that $\delta^2=d$. Note that 
\[
A_2B_2=D_4=(X^2-dY^2)(\mu^2X-dY^2),
\]
and $\gcd(A_1,A_2)=1$, $\gcd(B_1,B_2)=1$ because of \eqref{gcd=1}. Therefore, we have 
\begin{align*}
&A_2=\epsilon(X\pm\delta Y)(\mu X+\delta Y),\cr
&B_2= \epsilon^{-1}(X\mp\delta Y)(\mu X-\delta Y).
\end{align*}
where $\epsilon\in{\overline\f_q}\!^*$. It follows that 
\begin{align*}
D_3\,&=A_1B_2+A_2B_1\cr
&=\epsilon^{-1}(X\mp \delta Y)(\mu X-\delta Y)^2+\epsilon(X\pm\delta Y)(\mu X+\delta Y)^2\cr
&=(\epsilon^{-1}+\epsilon)\mu^2X^3+(\epsilon^{-1}-\epsilon)(\mp\delta^3)Y^3+\text{other terms}.
\end{align*}
Therefore,
\[
\begin{cases}
\epsilon^{-1}+\epsilon=4,\cr
\epsilon^{-1}-\epsilon=0.
\end{cases}
\]
In the above, the second equation gives $\epsilon=\pm 1$. Then $\epsilon^{-1}+\epsilon=2$ or $-2$, which are not equal to 4, which is a contradiction. Hence the claim is proved.

Now, by the Hasse-Weil bound,
\[
|V(D)|\ge q+1-(4-1)(4-2)q^{1/2}=q-6q^{1/2}+1.
\]
(Note that the curve $D(X,Y)=0$ has no rational points at $\infty$.) 

Given $(a_1,u)\in V(D)$, we may recover $a_0,b_0,b_1,s,t$ using  \eqref{du2}, \eqref{b0b1}, \eqref{s=t=}. In doing so, we need to require 
\begin{equation}\label{need}
u\ne0,\quad 2a_0+a_1\ne 0,\quad t(1+s+t)\ne 0.
\end{equation}
Under the condition $2a_0+a_1\ne 0$, we have 
\[
t(1+s+t)=\frac{a_0^2 a_1 (2+a_1) (1+a_0+a_1)^2}{(2 a_0+a_1)^2}.
\]
Note that $a_0\ne0$ and $1+a_0+a_1\ne 0$ since $X^2+a_1X+a_0$ is irreducible over $\f_q$. Therefore, the conditions in \eqref{need} are equivalent to
\[
u\ne0,\quad 2a_0+a_1\ne 0,\quad a_1(2+a_1)\ne 0,
\]
which are equivalent to, in terms of $a_1$ and $u$, 
\[
E(a_1,u)\ne 0,
\]
where 
\[
E(X,Y)=Y(2X+X^2-dY^2) X(2+X).
\]
By B\'ezout's theorem,
\[
|V(D)\cap V(E)|\le 4\cdot 5=20.
\]
Hence 
\[
|V(D)\setminus V(E)|=|V(D)|-|V(D)\cap V(E)|\ge q-6q^{1/2}-19>0.
\]
(Recall that we assumed $q\ge 71$.)

Now choose $(a_1,u)\in V(D)\setminus V(E)$, and let $a_0,b_0,b_1,s,t$ be given by \eqref{du2}, \eqref{b0b1} and \eqref{s=t=}. Then
\[
\theta(f_{s,t})=\frac{s^3}{t(1+s+t)}=27+a^2
\]
and 
\[
G_{s,t}=(X^2+a_1X+a_0)(X^2+b_1X+b_0),
\]
where $X^2+a_1X+a_0$ and $X^2+b_1X+b_0$ are irreducible over $\f_q$. Moreover, $\mu^2\ne 9$ implies that $X^2+a_1X+a_0\ne X^2+b_1X+b_0$. Therefore, $\alpha=a^2\in\Theta_{\rm b}-27$.  
\end{proof}

\end{document}
